\section{声子：从晶格振动到独立振子}

Gross 第 12 章从离子位移的二次型哈密顿量出发，经简正坐标对角化为独立振子，其量子即声子。

\subsection{谐振近似}

离子相对平衡位 $\mathbf{R}_\ell^0$ 的位移 $\mathbf{u}_\ell$，势能展到二次：
\begin{equation}
  V_{\mathrm{ion}}
  =\frac12\sum_{\ell\ell',\alpha\beta}
  u_{\ell\alpha}
  \Phi_{\alpha\beta}(\ell-\ell')
  u_{\ell'\beta}.
\end{equation}
力常数
\begin{equation}
  \Phi_{\alpha\beta}(\ell-\ell')
  =\frac{\partial^2 W}{\partial R_{\ell\alpha}\partial R_{\ell'\beta}}\Big|_{\mathrm{eq}}
\end{equation}
满足平移齐次性 $\sum_{\ell'}\Phi(\ell-\ell')=0$（整体平移不费力）。动能
$T=\sum_{\ell}M|\dot{\mathbf{u}}_\ell|^2/2$。

\subsection{Bloch 变换与动力学矩阵}

利用晶格周期边界，引入
\begin{equation}
  \mathbf{u}_\ell
  =\frac{1}{\sqrt{N}}
  \sum_{\mathbf{q}\in\mathrm{BZ}}
  e^{i\mathbf{q}\cdot\mathbf{R}_\ell^0}
  \mathbf{U}_{\mathbf{q}},
\end{equation}
势能化为
\begin{equation}
  V_{\mathrm{ion}}
  =\frac12\sum_{\mathbf{q}}
  \mathbf{U}_{\mathbf{q}}^\dagger
  A(\mathbf{q})
  \mathbf{U}_{\mathbf{q}},
\end{equation}
其中 $3\times3$ 矩阵
\begin{equation}
  A_{\alpha\beta}(\mathbf{q})
  =\sum_{\boldsymbol{\ell}}
  e^{-i\mathbf{q}\cdot\mathbf{R}_\ell^0}
  \Phi_{\alpha\beta}(\boldsymbol{\ell}).
\end{equation}
由 $\Phi_{\alpha\beta}(\boldsymbol{\ell})=\Phi_{\beta\alpha}(-\boldsymbol{\ell})$ 可证 $A(\mathbf{q})=A(\mathbf{q})^\dagger$ 且实对称（适当基下），故可被正交矩阵 $Y_{\mathbf{q}}$ 对角化：
\begin{equation}
  Y_{\mathbf{q}}^\dagger A(\mathbf{q})Y_{\mathbf{q}}
  =\mathrm{diag}(M\omega_{\mathbf{q}1}^2,M\omega_{\mathbf{q}2}^2,M\omega_{\mathbf{q}3}^2).
\end{equation}
令 $\mathbf{U}_{\mathbf{q}}=Y_{\mathbf{q}}\mathbf{Q}_{\mathbf{q}}$。动能在正交变换下不变（标量积保持），于是
\begin{equation}
  H_{\mathrm{ph}}
  =\sum_{\mathbf{q}\mu}
  \left(
  \frac{|P_{\mathbf{q}\mu}|^2}{2M}
  +\frac12 M\omega_{\mathbf{q}\mu}^2 |Q_{\mathbf{q}\mu}|^2
  \right).
\label{eq:Hph_classical}
\end{equation}

\subsection{色散的对称性}

由 $A(\mathbf{q})=A(-\mathbf{q})$（实力常数）得 $\omega_{\mathbf{q}\mu}=\omega_{-\mathbf{q}\mu}$；由晶格周期得 $\omega_{\mathbf{q}+\mathbf{G},\mu}=\omega_{\mathbf{q}\mu}$。故只需在第一 Brillouin 区计算。单原子原胞：三支声学模，长波 $\omega\sim c|\mathbf{q}|$（声速 $c=\partial\omega/\partial q|_{q=0}$）。多原子原胞另有光学支，$\omega(\mathbf{q}\to0)\to$ 有限。

\subsection{量子化}

\begin{equation}
  Q_{\mathbf{q}\mu}
  =\sqrt{\frac{1}{2M\omega_{\mathbf{q}\mu}}}
  \bigl(a_{\mathbf{q}\mu}+a_{-\mathbf{q}\mu}^\dagger\bigr),
  \qquad
  P_{\mathbf{q}\mu}
  =-i\sqrt{\frac{M\omega_{\mathbf{q}\mu}}{2}}
  \bigl(a_{\mathbf{q}\mu}-a_{-\mathbf{q}\mu}^\dagger\bigr),
\end{equation}
得
\begin{equation}
  H_{\mathrm{ph}}
  =\sum_{\mathbf{q}\mu}
  \omega_{\mathbf{q}\mu}
  \bigl(a_{\mathbf{q}\mu}^\dagger a_{\mathbf{q}\mu}+\tfrac12\bigr).
\end{equation}
三次及以上非谐项 $\propto u^3,u^4$ 引起声子--声子散射、热膨胀与有限热导，谐振近似下忽略。

\begin{note}
电子--离子耦合线性于 $\mathbf{u}$ 时给出 Fr\"ohlich 哈密顿量 \eqref{eq:frohlich}。下一章用二阶微扰导出声子介导的有效电子--电子相互作用。
\end{note}
